\documentclass[10pt,a4paper]{amsart}
\usepackage[margin=19mm]{geometry}
\usepackage{amsmath,amssymb,mathtools}
\usepackage[scr=rsfs,cal=euler]{mathalfa}
\usepackage{xfrac}
\usepackage[unicode,hidelinks,bookmarks=false]{hyperref}
\newcommand{\ie}{i.e.\ }
\newcommand{\cf}{cf.\ }
\newcommand{\resp}{respectively\ }
\newcommand{\Subsection}{\subsection}
\providecommand{\nobreakdash}{\nobreak\hbox{-}}
\setlength{\emergencystretch}{2em}
\multlinegap0pt
\usepackage{amssymb}
\usepackage{mathtools}
\usepackage{enumerate}
\newtheorem{lemma}{Lemma}
\newcommand{\B}{\mathbb{B}}
\newcommand{\Sn}{\mathbb{S}^{2n-1}}
\title{Contraction properties for holomorphic functions via isoperimetric stability on the Bergman ball}
\author{David Kalaj, Jian-Feng Zhu}
\date{Source: arXiv:2603.22524v2 (CC BY 4.0)}
\begin{document}
\maketitle

\noindent\emph{Source and licence.} Contraction properties for holomorphic functions via isoperimetric stability on the Bergman ball. Version arXiv:2603.22524v2 (CC BY 4.0), distributed by arXiv under the Creative Commons Attribution 4.0 International licence, http://creativecommons.org/licenses/by/4.0/, as recorded in the arXiv licence metadata for this exact version and shown on the abstract page https://arxiv.org/abs/2603.22524v2. Full licence notice: \texttt{licenses/arxiv-2603.22524-CC-BY-4.0.txt}.\par\smallskip

\noindent\textit{Editorial scope.} Counted unit: the article's lemma lem:IFT\_radial\_graph (source0.tex lines 1157--1178). The article's complete original proof of it is delivered with it (source0.tex lines 1180--1290, one \texttt{proof} environment of 111 lines). No mathematics is written, repaired or re-derived: the statement and the proof are the article's own lines, in the article's own notation, and no retained line is substituted or re-flowed.  Retained uncounted as the source-local context the proof needs: lines 1151--1153.  Nothing was omitted from any retained span, so the package carries no omission record.  No retained line contains a citation, so the wrapper carries no bibliography and no bibliography entry.  No retained line contains a figure environment, an image-inclusion command or drawing code, so no image is delivered and the package declares no asset dependency.  Editorial formatting only, in addition: an AMS \texttt{amsart} preamble that restates the article's own declarations and macros used by the retained lines, namely the environments \texttt{lemma} on separate counters, together with the standard packages those lines need.  The retained environments print in this document as Lemma 1 in order of appearance, whereas the article prints the same items with the numbering of its own full document.  The complete native source of the article as distributed in the arXiv e-print for this exact version is bundled unchanged as \texttt{sources/197023/source0.tex}.
\begin{equation}\label{eq:tail_condition}
\sup_{|z|\ge\rho}|f(z)|^p(1-|z|^2)^\alpha<t_-.
\end{equation}

\begin{lemma}[Radial graph lemma]\label{lem:IFT_radial_graph}
There exists a constant
\[
\delta_{\mathrm{rg}}=\delta_{\mathrm{rg}}(n,p,\alpha,\rho,\rho_-,\rho_+)>0
\]
such that if
\[
\|\phi\|_{C^1(\overline{\B_\rho})}\le \delta_{\mathrm{rg}},
\]
and \eqref{eq:tail_condition} holds, then for every \(t\in[t_-,t_+]\) the
level set \(\{u=t\}\) is a Lipschitz radial graph:
\begin{equation}\label{eq:rho_graph}
\{u=t\}=\{\rho_t(\omega)\,\omega:\ \omega\in\Sn\},
\end{equation}
where \(\rho_t:\Sn\to(\rho_-/2,(\rho_++\rho)/2)\) satisfies
\begin{equation}\label{eq:rho_est}
\|\rho_t-\rho_0(t)\|_{W^{1,\infty}(\Sn)}
\le
C\,\|\phi\|_{C^1(\overline{\B_\rho})},
\end{equation}
with \(C=C(n,p,\alpha,\rho,\rho_-,\rho_+)\).
\end{lemma}

\begin{proof}
Fix \(t\in[t_-,t_+]\) and \(\omega\in\Sn\). For \(\rho\in(0,1)\), set
\[
F(\rho,\omega):=\log u(\rho\omega)-\log t
=
p\log|1+\phi(\rho\omega)|+\alpha\log(1-\rho^2)-\log t.
\]
Choose \(\delta_{\mathrm{rg}}\) so small that
\[
\|\phi\|_{C^0(\overline{\B_\rho})}\le \frac12.
\]
Then \(|1+\phi|\ge \frac12\) on \(\overline{\B_\rho}\), so \(\log|1+\phi|\) is well-defined and smooth.

For \(\phi\equiv 0\), the equation \(F(\rho,\omega)=0\) is independent of \(\omega\) and has the unique solution
\(\rho=\rho_0(t)\in[\rho_-,\rho_+]\).

Differentiating in \(\rho\):
\[
\partial_\rho F(\rho,\omega)
=
p\,\partial_\rho\bigl(\log|1+\phi(\rho\omega)|\bigr)
-\frac{2\alpha\rho}{1-\rho^2}.
\]
Let
\[
I:=[\rho_-/2,(\rho_++\rho)/2].
\]
For $s\in I$, the second term is bounded away from $0$, since the function
$s\mapsto s/(1-s^2)$ is increasing on $(0,1)$ and hence
\[
\frac{2\alpha s}{1-s^2}
\ge
\frac{\alpha\rho_-}{1-\rho_-^2/4}
=:2c_0.
\]
On the other hand,
\[
\bigl|\partial_\rho\log|1+\phi(\rho\omega)|\bigr|
\le
\frac{|\nabla\phi(\rho\omega)\cdot\omega|}{|1+\phi(\rho\omega)|}
\le
2\|\nabla\phi\|_{L^\infty(\overline{\B_\rho})}
\le
2\|\phi\|_{C^1(\overline{\B_\rho})}.
\]
Hence, after shrinking \(\delta_{\mathrm{rg}}\) if necessary so that \(2p\,\delta_{\mathrm{rg}}\le c_0\), we obtain
\begin{equation}\label{eq:drF_neg_appendix}
\partial_\rho F(\rho,\omega)\le -c_0<0
\qquad\text{for all }(\rho,\omega)\in I\times\Sn.
\end{equation}
Thus, for each \(\omega\), the map \(\rho\mapsto F(\rho,\omega)\) is strictly decreasing on \(I\).

Since
\[
F(\rho_0(t),\omega)=p\log|1+\phi(\rho_0(t)\omega)|,
\]
we have
\[
|F(\rho_0(t),\omega)|
\le
C\,\|\phi\|_{C^0(\overline{\B_\rho})}.
\]
After decreasing $\delta_{\mathrm{rg}}$, the values of $F$ at the endpoints of
$I$ have opposite signs, uniformly in $t$ and $\omega$.  Thus
\eqref{eq:drF_neg_appendix} and the intermediate value theorem give a unique
\(\rho_t(\omega)\in I\) such that \(F(\rho_t(\omega),\omega)=0\), and moreover
\begin{equation}\label{eq:rho_Linf_est_appendix}
|\rho_t(\omega)-\rho_0(t)|
\le
C\,\|\phi\|_{C^0(\overline{\B_\rho})}.
\end{equation}
It remains to exclude roots outside $I$. For the model function, compactness
of $[t_-,t_+]$ gives a uniform positive gap on
$[0,\rho_-/2]$ and a uniform negative gap on
$[(\rho_++\rho)/2,\rho]$. The perturbation term
$p\log|1+\phi|$ is uniformly $O(\delta_{\mathrm{rg}})$; after one final
decrease of $\delta_{\mathrm{rg}}$, these signs persist. The tail condition
\eqref{eq:tail_condition} excludes roots for $|z|\ge\rho$. Hence the root
found in $I$ is the only root on each ray, and this proves
\eqref{eq:rho_graph} for the entire level set.

Differentiating \(F(\rho_t(\omega),\omega)=0\) tangentially on \(\Sn\) yields
\[
\nabla_\omega \rho_t(\omega)
=
-\frac{\nabla_\omega F(\rho_t(\omega),\omega)}{\partial_\rho F(\rho_t(\omega),\omega)}.
\]
Since \(\log(1-\rho^2)\) is independent of \(\omega\),
\[
\nabla_\omega F(\rho,\omega)
=
p\,\nabla_\omega\bigl(\log|1+\phi(\rho\omega)|\bigr).
\]
Using \(|1+\phi|\ge \frac12\),
\[
\bigl|\nabla_\omega\log|1+\phi(\rho\omega)|\bigr|
\le
C\,|\nabla_\omega(\phi(\rho\omega))|
\le
C\,\rho\,\|\nabla\phi\|_{L^\infty(\overline{\B_\rho})}
\le
C\,\|\phi\|_{C^1(\overline{\B_\rho})}.
\]
Together with \eqref{eq:drF_neg_appendix}, this gives
\begin{equation}\label{eq:rho_W1inf_est_appendix}
\|\nabla_\omega \rho_t\|_{L^\infty(\Sn)}
\le
C\,\|\phi\|_{C^1(\overline{\B_\rho})}.
\end{equation}
Combining \eqref{eq:rho_Linf_est_appendix} and \eqref{eq:rho_W1inf_est_appendix} proves \eqref{eq:rho_est}.
\end{proof}

\end{document}
